% Source: Evans, "Partial Differential Equations" (2nd ed., AMS Graduate Studies in Mathematics vol. 19, 2010),
% Chapter A "Notation", PDF pages 615-622 of MathBooks/Evans Partial Differential Equations(1).pdf.
% Transcribed page-by-page by vision subagents, then merged and restructured into
% leanblueprint conventions (semantic \label, bracketed titles, \uses dependency edges) below.


\section{Basic Notation}
\label{sec:basic-notation}

\subsection{Notation for Matrices}
\label{sec:notation-for-matrices}

\begin{definition}[Matrix Notation]
\label{def:matrix-notation-basic}
\dcref{appA:A.1-def-1}
\group{Matrix Notation}
(i) We write $A = ((a_{ij}))$ to mean $A$ is an $m \times n$ matrix with $(i,j)^{th}$ entry $a_{ij}$. (Sometimes, as in \S8.1.4, it will be convenient to use superscripts to denote rows.)

A diagonal matrix is denoted $\text{diag}(d_1, \ldots, d_n)$.

(ii) $M^{m \times n}$ = space of real $m \times n$ matrices.

$\mathbb{S}^{n \times n}$ = space of real symmetric $n \times n$ matrices.

(iii) $\text{tr}\, A$ = trace of the matrix $A$.

(iv) $\det A$ = determinant of the matrix $A$.

(v) $\text{cof}\, A$ = cofactor matrix of $A$ (See \S8.1.4).

(vi) $A^T$ = transpose of the matrix $A$.

(vii) If $A = ((a_{ij}))$ and $B = ((b_{ij}))$ are $m \times n$ matrices, then
$$A : B = \sum_{i=1}^{m} \sum_{j=1}^{n} a_{ij} b_{ij},$$
and the Frobenius norm of $A$ is
$$|A| = (A : A)^{1/2} = \left(\sum_{i=1}^{m} \sum_{j=1}^{n} a_{ij}^2\right)^{1/2}.$$
\end{definition}

\begin{definition}[Quadratic Forms and Related Matrix Notation]
\label{def:matrix-quadratic-form-notation}
\dcref{appA:A.1}
\group{Matrix Notation}
(viii) If $A \in \mathbb{S}^{n \times n}$ and, as below, $x = (x_1, \ldots, x_n) \in \R^n$, the corresponding quadratic form is $x \cdot Ax = \sum_{i,j=1}^{n} a_{ij} x_i x_j$.

(ix) If $A \in \mathbb{S}^{n \times n}$, we write
$$A \geq \theta I$$
if $x \cdot Ax \geq \theta |x|^2$ for all $x \in \R^n$.

(x) We sometimes will write $yA$ to mean $A^T y$, for $A \in M^{m \times n}$ and $y \in \R^m$.
\end{definition}

\section{Geometric Notation}
\label{sec:geometric-notation}

\begin{definition}[Geometric Notation]
\label{def:geometric-notation}
\dcref{appA:A.2}
\group{Geometric Notation}
(i) $\R^n = n$-dimensional real Euclidean space, $\R = \R^1$.

(ii) $e_i = (0, \ldots, 0, 1, \ldots, 0) = i^{\text{th}}$ standard coordinate vector.

(iii) A typical point in $\R^n$ is $x = (x_1, \ldots, x_n)$. We will also, depending upon the context, regard $x$ as a row or column vector.

(iv) $\R_+^n = \{x = (x_1, \ldots, x_n) \in \R^n \mid x_n > 0\} = $ open upper half-space. $\R_+ = \{x \in \R \mid x > 0\}$.

(v) A typical point in $\R^{n+1}$ will often be denoted as $(x,t) = (x_1, \ldots, x_n, t)$, and we usually interpret $t = x_{n+1} = $ time. A point $x \in \R^n$ will sometimes be written $x = (x', x_n)$ for $x' = (x_1, \ldots, x_{n-1}) \in \R^{n-1}$.

(vi) $U$, $V$, and $W$ usually denote open subsets of $\R^n$. We write
$$V \subset\subset U$$
if $V \subset \bar{V} \subset U$ and $\bar{V}$ is compact, and say $V$ is \textit{compactly contained} in $U$.

(vii) $\partial U = $ boundary of $U$, $\bar{U} = U \cup \partial U = $ closure of $U$.

(viii) $U_T = U \times (0, T]$.

(ix) $\Gamma_T = \bar{U}_T - U_T = $ parabolic boundary of $U_T$.

(x) $B^o(x,r) = \{y \in \R^n \mid |x - y| < r\} = $ open ball in $\R^n$ with center $x$ and radius $r > 0$.

(xi) $B(x, r) = \text{closed ball with center } x\text{, radius } r > 0$.

(xii) $C(x, t, r) = \{y \in \R^n, s \in \R \mid |x - y| \le r, t - r^2 \le s \le t\} = \text{closed cylinder with top center } (x, t)\text{, radius } r > 0\text{, height } r^2$.

(xiii) $\alpha(n) = \text{volume of unit ball } B(0, 1) \text{ in } \R^n = \dfrac{\pi^{n/2}}{\Gamma(\frac{n}{2}+1)}$,
$n\alpha(n) = \text{surface area of unit sphere } \partial B(0, 1) \text{ in } \R^n$.

(xiv) If $a = (a_1, \ldots, a_n)$ and $b = (b_1, \ldots, b_n)$ belong to $\R^n$,
$$a \cdot b = \sum_{i=1}^n a_i b_i, \quad |a| = \left(\sum_{i=1}^n a_i^2\right)^{1/2}.$$

(xv) $\C^n = n\text{-dimensional complex space, } \C = \text{complex plane}$.
If $z \in \C$, we write $\text{Re}(z)$ for the real part of $z$, and $\text{Im}(z)$ for the imaginary part.
\end{definition}

\section{Notation for Functions}
\label{sec:notation-for-functions}

\begin{definition}[Notation for Functions]
\label{def:function-notation-basic}
\dcref{appA:A.3}
\group{Function Notation}
\uses{def:geometric-notation}
(i) If $u : U \to \R$, we write
$$u(x) = u(x_1, \ldots, x_n) \quad (x \in U).$$
We say $u$ is \textit{smooth} provided $u$ is infinitely differentiable.

(ii) If $u, v$ are two functions, we write
$$u \equiv v$$
to mean that $u$ is identically equal to $v$; that is, the functions $u, v$ agree for all values of their arguments.

Let us also set
$$u := v$$
to define $u$ as equaling $v$.

The support of a function $u$ is denoted
$$\spt u.$$

(iii) $u^+ = \max(u, 0)$, $u^- = -\min(u, 0)$, $u = u^+ - u^-$, $|u| = u^+ + u^-$.

The \textit{sign} function is
$$\text{sgn}(x) = \begin{cases} 1 & \text{if } x > 0 \\ 0 & \text{if } x = 0 \\ -1 & \text{if } x < 0. \end{cases}$$

(iv) If $\mathbf{u} : U \to \R^m$, we write
$$\mathbf{u}(x) = (u^1(x), \ldots, u^m(x)) \quad (x \in U).$$
The function $u^k$ is the $k^{\text{th}}$ \textit{component} of $\mathbf{u}$, $k = 1, \ldots, m$.

(v) If $\Sigma$ is a \textit{smooth} $(n-1)$-dimensional \textit{surface} in $\R^n$, we write
$$\int_\Sigma f \, dS$$
for the integral of $f$ over $\Sigma$, with respect to $(n-1)$-dimensional \textit{surface} measure. If $C$ is a curve in $\R^n$, we denote by
$$\int_C f \, dl$$
the integral of $f$ over $C$ with respect to arclength.

(vi) \textit{Averages}:
$$\fint_{B(x,r)} f \, dy = \frac{1}{\alpha(n)r^n} \int_{B(x,r)} f \, dy = \text{average of } f \text{ over the ball } B(x,r)$$
and
$$\fint_{\partial B(x,r)} f \, dS = \frac{1}{n\alpha(n)r^{n-1}} \int_{\partial B(x,r)} f \, dS$$
$$= \text{average of } f \text{ over the sphere } \partial B(x,r).$$

(vii) $\chi_E(x) = \begin{cases} 1 & \text{if } x \in E \\ 0 & \text{if } x \notin E \end{cases}$ $\chi_E$ is the \textit{indicator function} of $E$.

(viii) A function $u : U \to \R$ is called \textit{Lipschitz continuous} if
$$|u(x) - u(y)| \leq C|x - y|$$
for some constant $C$ and all $x, y \in U$. We write
$$\text{Lip}[u] := \sup_{\substack{x,y \in U \\ x \neq y}} \frac{|u(x) - u(y)|}{|x - y|}.$$

(ix) The convolution of the functions $f, g$ is denoted
$$f * g.$$
\end{definition}

\subsection{Notation for Derivatives}
\label{sec:notation-for-derivatives}

\begin{definition}[Notation for Derivatives]
\label{def:derivative-notation-multiindex}
\dcref{appA:A.3-derivatives}
\group{Derivative Notation}
Assume $u : U \to \R$, $x \in U$.

(i) $\pd{u}{x_i}(x) = \lim\limits_{h \to 0} \dfrac{u(x+he_i)-u(x)}{h}$, provided this limit exists.

(ii) We usually write $u_{x_i}$ for $\pd{u}{x_i}$.

(iii) Similarly, $\dfrac{\partial^2 u}{\partial x_i \partial x_j} = u_{x_i x_j}$, $\dfrac{\partial^3 u}{\partial x_i \partial x_j \partial x_k} = u_{x_i x_j x_k}$, etc.

(iv) \textit{Multiindex Notation}:

(a) A vector of the form $\alpha = (\alpha_1, \ldots, \alpha_n)$, where each component $\alpha_i$ is a nonnegative integer, is called a \textit{multiindex of order}
$$|\alpha| = \alpha_1 + \cdots + \alpha_n.$$

(b) Given a multiindex $\alpha$, define
$$D^\alpha u(x) := \frac{\partial^{|\alpha|} u(x)}{\partial x_1^{\alpha_1} \cdots \partial x_n^{\alpha_n}} = \partial_1^{\alpha_1} \cdots \partial_n^{\alpha_n} u.$$

(c) If $k$ is a nonnegative integer,
$$D^k u(x) := \{D^\alpha u(x) \mid |\alpha| = k\},$$
the set of all partial derivatives of order $k$. Assigning some ordering to the various partial derivatives, we can also regard $D^k u(x)$ as a point in $\R^{n^k}$.

(d)
$$|D^k u| = \left(\sum_{|\alpha|=k} |D^\alpha u|^2\right)^{1/2}.$$

(e) \textit{Special Cases}: If $k = 1$, we regard the elements of $Du$ as being arranged in a vector:
$$Du = (u_{x_1}, \ldots, u_{x_n}) = \textit{gradient vector}.$$

If $k = 2$, we regard the elements of $D^2 u$ as being arranged in a matrix:
$$D^2 u = \begin{pmatrix}
\dfrac{\partial^2 u}{\partial x_1^2} & \cdots & \dfrac{\partial^2 u}{\partial x_1 \partial x_n} \\[1em]
\vdots & \ddots & \vdots \\[0.5em]
\dfrac{\partial^2 u}{\partial x_n \partial x_1} & \cdots & \dfrac{\partial^2 u}{\partial x_n^2}
\end{pmatrix}_{n \times n} = \textit{Hessian matrix}.$$

(v)
$$\Delta u = \sum_{i=1}^n u_{x_i x_i} = \text{tr}(D^2 u) = \textit{Laplacian of } u.$$

(vi) We sometimes employ a subscript attached to the symbols $D$, $D^2$, etc. to denote the variables being differentiated. For example if $u = u(x,y)$ ($x \in \R^n$, $y \in \R^m$), then $D_x u = (u_{x_1}, \ldots, u_{x_n})$, $D_y u = (u_{y_1}, \ldots, u_{y_m})$.
\end{definition}

\subsection{Function Spaces}
\label{sec:function-spaces-notation}

\begin{definition}[Function Spaces]
\label{def:function-space-notation}
\dcref{appA:A.3-function-spaces}
\group{Function Spaces}
\uses{def:derivative-notation-multiindex,def:geometric-notation}
(i) $C(U) = \{u : U \to \R \mid u \text{ continuous}\}$

$\overline{C}(U) = \{u \in C(U) \mid u \text{ uniformly continuous}\}$

$C^k(U) = \{u : U \to \R \mid u \text{ is } k\text{-times continuously differentiable}\}$

$\overline{C^k(U)} = \{u \in C^k(U) \mid D^\alpha u \text{ is uniformly continuous for all } |\alpha| \leq k\}$.

Thus if $u \in \overline{C^k(U)}$, then $D^\alpha u$ continuously extends to $\overline{U}$ for each multiindex $\alpha$, $|\alpha| \leq k$.

(ii) $C^\infty(U) = \{u : U \to \R \mid u \text{ is infinitely differentiable}\} = \bigcap_{k=0}^{\infty} C^k(U)$

$\overline{C^\infty(U)} = \bigcap_{k=0}^{\infty} \overline{C^k(U)}$.

(iii) $C_c(U)$, $C_c^k(U)$, etc. denote these functions in $C(U)$, $C^k(U)$, etc. with \textit{compact support}.

(iv) $L^p(U) = \{u : U \to \R \mid u \text{ is Lebesgue measurable, } \|u\|_{L^p(U)} < \infty\}$, where
$$\|u\|_{L^p(U)} = \left(\int_U |u|^p dx\right)^{1/p} \quad (1 \leq p < \infty).$$
$L^\infty(U) = \{u : U \to \R \mid u \text{ is Lebesgue measurable, } \|u\|_{L^\infty(U)} < \infty\}$, where
$$\|u\|_{L^\infty(U)} = \text{ess sup}_U |u|.$$
$L^p_{\text{loc}}(U) = \{u : U \to \R \mid u \in L^p(V) \text{ for each } V \subset\subset U\}$.

(See also \S D.1.)

(v) $\|Du\|_{L^p(U)} = \||Du|\|_{L^p(U)}$

$\|D^2 u\|_{L^p(U)} = \||D^2 u|\|_{L^p(U)}$.

(vi) $W^{k,p}(U)$, $H^k(U)$, etc. $(k = 0, 1, 2, \ldots, 1 \leq p \leq \infty)$ denote Sobolev spaces: see Chapter 5.

(vii) $C^{k,\beta}(U)$, $C^{k,\beta}(\overline{U})$ $(k = 0, \ldots, 0 < \beta \leq 1)$ denote H\"older spaces: see Chapter 5.

(viii) \textit{Functions of $x$ and $t$.} It is occasionally useful to introduce spaces of functions with differing smoothness in the $x$- and $t$-variables, although there is no standard notation for such spaces. We will for this book write
$$C_T^2(U_T) = \{u : U_T \to \R \mid u, D_x u, D_x^2 u, u_t \in C(U_T)\}.$$
In particular, if $u \in C_T^2(U_T)$, then $u, D_x u$, etc. are continuous up to the top $U \times \{t = T\}$.
\end{definition}

\section{Vector-Valued Functions}
\label{sec:vector-valued-functions}

\begin{definition}[Vector-Valued Functions]
\label{def:vector-valued-function-notation}
\dcref{appA:A.4}
\group{Vector-Valued Functions}
\uses{def:derivative-notation-multiindex,def:function-notation-basic}
(i) If now $m > 1$ and $\mathbf{u} : U \to \R^m$, $\mathbf{u} = (u^1, \ldots, u^m)$, we define
$$D^\alpha \mathbf{u} = (D^\alpha u^1, \ldots, D^\alpha u^m) \quad \text{for each multiindex } \alpha.$$
Then
$$D^k \mathbf{u} = \{D^\alpha \mathbf{u} \mid |\alpha| = k\}$$
and
$$|D^k \mathbf{u}| = \left(\sum_{|\alpha|=k} |D^\alpha \mathbf{u}|^2\right)^{1/2},$$
as before.

(ii) In the special case $k = 1$, we write
$$D\mathbf{u} = \begin{pmatrix} \dfrac{\partial u^1}{\partial x_1} & \cdots & \dfrac{\partial u^1}{\partial x_n} \\[1em] \vdots & \ddots & \vdots \\[0.5em] \dfrac{\partial u^m}{\partial x_1} & \cdots & \dfrac{\partial u^m}{\partial x_n} \end{pmatrix}_{m \times n} = \text{gradient matrix}.$$

(iii) If $m = n$, we have
$$\Div \mathbf{u} = \text{tr}(D\mathbf{u}) = \sum_{i=1}^{n} u^i_{x_i} = \text{divergence of } \mathbf{u}.$$

(iv) The spaces $C(U;\R^m)$, $L^p(U;\R^m)$, etc. consist of those functions $\mathbf{u} : U \to \R^m$, $\mathbf{u} = (u^1, \ldots, u^m)$, with $u^i \in C(U)$, $L^p(U)$, etc. $(i = 1, \ldots, m)$.
\end{definition}

\begin{remark}[Sub- and Superscripts for Boldface Mappings]
\label{rem:boldface-convention-sub-superscripts}
\dcref{appA:A.4-remark}
\group{Notation Conventions}
\uses{def:vector-valued-function-notation}
As illustrated above, we will adhere to the convention of setting in boldface \textit{mappings} which take values in $\R^m$ for $m > 1$ (or else in Banach or Hilbert spaces). The component functions of such mappings will be given \textit{superscripts}. On the other hand, a typical \textit{point} $x \in \R^n$ is not boldface and has components with \textit{subscripts}, $x = (x_1, \ldots, x_n)$.

Matrix-valued mappings will also be set in boldface, and their component functions written with either superscripts or a mixture of sub- and superscripts, depending on the context. $\square$
\end{remark}

\section{Notation for Estimates}
\label{sec:notation-for-estimates}

\begin{remark}[Constants]
\label{rem:constant-notation-convention}
\dcref{appA:A.5-constants}
\group{Asymptotic Notation}
We employ the letter $C$ to denote any constant that can be explicitly computed in terms of known quantities. The exact value denoted by $C$ may therefore change from line to line in a given computation. The big advantage is that our calculations will be simpler looking, since we continually absorb ``extraneous'' factors into the term $C$.
\end{remark}

\begin{definition}[Big-Oh and Little-Oh Notation]
\label{def:big-oh-little-oh-notation}
\dcref{appA:A.5-definitions}
\group{Asymptotic Notation}
(i) (Big-oh notation.) \textit{We write}
$$f = O(g) \quad \text{as } x \to x_0,$$
\textit{provided there exists a constant } $C$ \textit{ such that}
$$|f(x)| \leq C|g(x)|$$
\textit{for all } $x$ \textit{ sufficiently close to } $x_0$.

(ii) (Little-oh notation.) \textit{We write}
$$f = o(g) \quad \text{as } x \to x_0,$$
\textit{provided}
$$\lim_{x \to x_0} \frac{|f(x)|}{|g(x)|} = 0.$$
\end{definition}

\begin{remark}[Implicit Limits in $O$/$o$ Notation]
\label{rem:big-oh-little-oh-implicit-limit}
\dcref{appA:A.5-remark}
\group{Asymptotic Notation}
\uses{def:big-oh-little-oh-notation}
The expression ``$O(g)$'' (or ``$o(g)$'') is not by itself defined. There must always be an accompanying limit, for example ``as $x \to x_0$'' above, although this limit is often implicit. $\square$
\end{remark}

\section{Some Comments about Notation}
\label{sec:comments-about-notation}

The foregoing notation is largely standard within the PDE literature, with a few significant exceptions:

\begin{remark}[Why ``$Du$'' Instead of ``$\nabla u$'']
\label{rem:gradient-notation-choice-du}
\dcref{appA:A.6-i}
\group{Notation Conventions}
\uses{def:derivative-notation-multiindex}
We employ the symbol ``$Du$'', and not ``$\nabla u$'', to denote the gradient of the function $u$. The reason is that ``$D^2u$'' then naturally denotes the Hessian matrix of $u$, whereas ``$\nabla^2 u$'' would be confused with the Laplacian. The multiindex notation also looks better with the letter $D$.
\end{remark}

\begin{remark}[Why ``$U$'' Instead of ``$\Omega$'']
\label{rem:domain-notation-choice-u}
\dcref{appA:A.6-ii}
\group{Notation Conventions}
\uses{def:geometric-notation}
Most books and papers on partial differential equations denote by ``$\Omega$'' the open subset of $\R^n$ within which a given PDE holds.

As indicated above, we will instead mostly use the symbol ``$U$'' for such a region. The advantages are several. First of all, since a typical solution is denoted $u$, it makes sense to denote its domain by $U$, and not to switch to a Greek letter. Furthermore, once we call a given open set $U$, the letters $V$ and $W$ are then available for subregions.

Lastly, it is important to save $\Omega$ as the standard symbol for a probability space. Many important partial differential equations have probabilistic representation formulas (cf.\ Freidlin \cite{Freidlin1985}), and although such are beyond the scope of this book, it seems wise to avoid the possibility of future notational confusion.
\end{remark}
